%%=============================================================================
%% Lijst van afkortingen
%%=============================================================================

\chapter*{\IfLanguageName{dutch}{Lijst van afkortingen}{List of abbreviations}}%
\label{ch:afkortingen}
\addcontentsline{toc}{chapter}{\IfLanguageName{dutch}{Lijst van afkortingen}{List of abbreviations}}

\begin{description}[leftmargin=3.5cm,labelwidth=3cm,style=nextline,font=\bfseries]
    \item[BMC]      BMC Software, leverancier van mainframe-managementsoftware en bron van het jaarlijkse \textit{Mainframe Survey} rapport.
    \item[CD]       Core Drive, één van de acht psychologische motivatoren binnen het Octalysis-framework van \textcite{Chou2015}.
    \item[CICS]     Customer Information Control System, IBM-transactieserver voor mainframe-applicaties.
    \item[COBOL]    Common Business Oriented Language, procedurele programmeertaal, oorspronkelijk uit 1960, nog steeds dominant op mainframes.
    \item[CRUD]     Create, Read, Update, Delete, standaardoperaties op een datastore.
    \item[DB2]      Database 2, relationele database van IBM, courant gekoppeld aan COBOL-applicaties op z/OS.
    \item[DV]       Deelvraag, in dit document genummerd DV1 t.e.m.\ DV5 (zie sectie~\ref{sec:onderzoeksvraag}).
    \item[ECB]      European Central Bank.
    \item[IBM]      International Business Machines, historische ontwikkelaar van het mainframe-platform en van \textit{Enterprise COBOL for z/OS}.
    \item[IDE]      Integrated Development Environment.
    \item[ING]      ING Group, internationale bank-/verzekeringsgroep, voorbeeld van een grote COBOL-gebruiker.
    \item[JCL]      Job Control Language, scriptingtaal voor het inplannen van batchjobs op IBM-mainframes.
    \item[MDA]      Mechanics, Dynamics, Aesthetics, formeel framework voor game-design van \textcite{Hunicke2004}.
    \item[OOP]      Object-Oriented Programming, object-georiënteerd programmeren.
    \item[PBL]      Points, Badges, Leaderboards, de meest bekende componenten van gamification \autocite{KevinWerbach2020}.
    \item[PDF]      Portable Document Format, in dit document ook de gebruikte vormgeving van de traditionele leerconditie (``PDF-pack'').
    \item[PIC]      \texttt{PICTURE}-clausule in COBOL, definieert opslagformaat en datatype van een variabele.
    \item[PoC]      Proof of Concept, tastbaar prototype dat de haalbaarheid van een idee demonstreert.
    \item[RAS]      Reliability, Availability, Serviceability, kernkwaliteitsattributen van mainframe-platformen \autocite{Ebbers2011}.
    \item[RLS]      Row-Level Security, rij-gebaseerde toegangscontrole, hier op niveau van Supabase/PostgreSQL.
    \item[RPG]      Report Program Generator, andere legacy-taal van IBM, vooral op AS/400 en IBM~i.
    \item[SMD]      Standardised Mean Difference, gestandaardiseerd gemiddeld verschil, een effectgroottemaat in meta-analyses.
    \item[SUS]      System Usability Scale, tien-item Likert-vragenlijst voor de bruikbaarheid van een systeem (vermeld als aanbeveling voor empirische opvolging).
    \item[TI]       Toegepaste Informatica, opleiding aan HOGENT waaruit deze bachelorproef voortkomt.
    \item[UI]       User Interface, gebruikersinterface.
    \item[z/OS]     IBM-besturingssysteem voor de mainframe-familie System z.
\end{description}
